\documentclass[../main.tex]{subfiles}
\begin{document}
%==========================================TIÊU ĐỀ TẬP========================================
\begin{table}[h]
    \begin{adjustwidth}{-1cm}{}
        \begin{tabular}{>{\centering\arraybackslash}p{6cm}|>{\raggedright\arraybackslash}p{12.5cm}}
            \multirow{5}{6cm}
            % Cột bên trái: ÉPISODE và số tập
            {\\[-40pt]
                \flaregothic\fontsize{20pt}{20pt}\selectfont \centering
                \color{eptype}ÉPISODE\color{black}\\
                \burbank\fontsize{72pt}{72pt}\selectfont
                \color{epnum}38\color{black}\\[6pt]
                \hrule
                \vspace{8pt}
                \flaregothic\fontsize{20pt}{20pt}\selectfont
                \color{eptype}SPÉCIAL\color{black}\\
                \burbank\fontsize{72pt}{72pt}\selectfont
                \color{epnum}VI\color{black}\\[6pt]
                \cafeteria\fontsize{20pt}{20pt}\selectfont\color{epnum}(S-6)\color{black}
                \\[18pt]
                \flaregothic\fontsize{18pt}{18pt}\selectfont
                \color{gray}COMPILATION\\
                \color{epattr}ÉPISODE PERDU
                \color{black}
            }
             &
            % Dòng thứ nhất: tên tập tiếng Nhật
            {\fotskip\fontsize{18pt}{18pt}\selectfont \ruby{劇}{げき}\ruby{劇場}{じょう}\ruby{版}{ばん}ホロライブ！
            }                                                                                      \\[6pt]
            % Dòng thứ hai: tên tập tiếng Anh
             & {\cafeteria\fontsize{18pt}{18pt}\selectfont \hll\dots\ the Movie! - The Prelude}                                                                             \\[4pt]
            % Dòng thứ ba: tên tập tiếng Việt
             & {\vnmsans\fontsize{18pt}{18pt}\selectfont Phim\dots\ nhưng không phải là phim!}                                                                                 \\[8pt]
            % Dòng thứ tư: Nhân vật xuất hiện
             & {\caviar{\fontsize{14pt}{14pt}\selectfont Nhân vật: \emp{kuon1} \emp{ryuuhou1} \emp{achan} \emp{sora} \emp{roboco1} \emp{miko2} \emp{mel} \emp{fubuki} \emp{matsuri} \emp{haato} \emp{aqua} \emp{shion} \emp{ayame} \emp{choco} \emp{subaru} \emp{okayu} \emp{korone} \emp{pekora} \emp{rushia} \emp{noel}}}
            \\[8pt]
            % Dòng cuối cùng: Thời gian diễn ra của tập (thường sẽ là ngày phát hành)
             & {\continuum\fontsize{16pt}{16pt}\selectfont Ngày phát hành: 26/01/2020}                                                                           \\[18pt]
        \end{tabular}
    \end{adjustwidth}
\end{table}
%=========================================TRUYỆN CHÍNH========================================
\color{black}

\txtbx[2]{{\color{vol} \uline{Tóm tắt:}} Khi Kuon vắng mặt để về quê ăn Tết, văn phòng \textit{hololive} rơi vào trạng thái lười biếng cực độ. Để xốc lại tinh thần, các cô gái nảy ra ý tưởng thực hiện một bộ phim điện ảnh hoành tráng. Tuy nhiên, những cuộc tranh cãi nảy lửa về kịch bản—từ đồ bơi mùa đông đến chiếm đóng đài truyền hình—đã biến dự án thành một mớ hỗn độn. Sora xuất hiện với thông báo chấn động rằng bộ phim chỉ dài hai phút rưỡi và sắp kết thúc, nhưng một tiếng nói bí ẩn đã mở ra một hành trình hoàn toàn mới về trải nghiệm Tết Etsunan.}

Một ngày cuối tháng Giêng năm 2020. Bầu không khí lạnh giá của mùa đông vẫn còn vương vấn trên những con phố Tokyo, nhưng bên trong văn phòng \hll, sự náo nhiệt thường thấy đã bị thay thế bởi một sự yên ắng đến rợn người.

Hikawa Kuon, linh hồn và cũng là ``cỗ máy vận hành'' của văn phòng, đã chính thức nghỉ phép để trở về Etsunan đón Tết Nguyên Đán. Cậu rời đi ngay sau buổi livestream lớn nhất trong tháng, để lại một mẩu giấy nhắn đầy vẻ đe dọa trên bảng tin: ``Anh đi vắng vài hôm, mọi người tự quản nhé. Nhớ đừng làm loạn, nếu không khi về anh sẽ trừ sạch lương tháng sau đấy.''

Ban đầu, các cô gái đồng ý với sự tự tin đầy hào hứng (có lẽ bởi vì họ không chịu sự quản lý của một cậu Tổng nghiêm túc), nhưng giờ đây, khi không có Kuon để tổ chức, điều hành, và đôi khi pha trò châm chọc, cả nhóm rơi vào trạng thái lười biếng đến kỳ lạ.

\img{3-9a}

Matsuri đứng thẫn thờ, thả lỏng người xuống đất. Cô khẽ rên rỉ, giọng điệu đầy vẻ chán chường: ''Ôi xời\dots\ Chán quá đi mất\dots\ Không có ai để trêu, cũng chẳng có ai mắng mình nữa.''

Ở góc phòng, Roboco đang cầm một bát thức ăn, lơ đãng cho chú chó của Mio\footnote{Xem lại {\flaregothic ÉPISODE 21}, quyển số Hai.} ăn. Chú chó vẫy đuôi hào hứng, hoàn toàn trái ngược với vẻ mặt uể oải của nàng Người máy. Fubuki ngồi ngay kế bên, chống cằm nhìn chăm chú vào một điểm vô định trên tường: ''Một ngày thật yên bình\dots\ yên bình đến mức mình sắp hóa thạch rồi đây.''

Nàng Cổ động khẽ ngẩng đầu lên, ánh mắt lờ đờ nhìn cô bạn cáo trắng: ''Bà có thấy lạ không? Hay là do anh Tổng không ở đây nên mọi thứ mới thành ra thế này?''

Nàng Cáo trắng thở dài, chậm rãi đáp lời: ''Tôi cũng không biết nữa\dots\ Nhưng quả thật là không có Kuon-senpai ở đây, bầu không khí nó cứ\dots\ khang khác thế nào ấy. Thiếu thiếu một cái gì đó rất khó tả.''

Cảnh tượng trong văn phòng lúc này chẳng khác gì một trạm xe buýt bỏ hoang vào lúc xế chiều. Mio, người duy nhất còn giữ được chút năng lượng, đứng giữa phòng và nhìn quanh với vẻ không hài lòng chút nào. Cô chống hai tay ngang hông, hét lớn: ''Này! Mấy người tỉnh lại đi! Đừng có tỏ ra lười biếng như thế nữa chứ!''

\img{3-9b}

Fubuki và Matsuri đồng loạt quay đầu lại, nhìn Mio bằng ánh mắt như thể cô vừa nói một ngôn ngữ xa lạ nào đó. Nàng Cáo trắng nhướng mày hỏi: ''Gì thế Mio?'' theo sau là cô bạn Cổ động cũng thêm vào, giọng uể oải: ``Sao tự nhiên lại sung thế?''

Mio nhắm mắt, hít một hơi thật sâu để kiềm chế cơn giận đang chực trào: ''Trời ạ, hai người này! Chúng ta đang ở trong một bộ phim điện ảnh đấy! Các người có hiểu không? Đây là cơ hội để chúng ta làm một cái gì đó thật hoành tráng, thật khác biệt để khán giả phải trầm trồ! Sao lại cứ thích ngồi không nhạt nhẽo thế này?''

Cặp đôi Thế hệ thứ Nhất ngơ ngác nhìn nhau, rồi đồng thanh thốt lên: ''Ể!? Phim điện ảnh á!?'' ``Chứ không phải mấy cái video ngắn mà chúng ta vẫn làm hàng tuần sao?''

Mio đập tay cái 'rầm' lên bàn, giọng điệu đầy quyết liệt: ''Không! Đây là MOVIE! Là phim chiếu rạp hẳn hoi đấy! Nào, tất cả đứng dậy ngay cho tôi! Chúng ta phải xây dựng kịch bản ngay lập tức!''

Ở một góc khác, Noel và Rushia đang đắm chìm trong ván cờ lật\footnote{Xem lại {\flaregothic ÉPISODE 31} trong quyển này.} đầy căng thẳng. Nữ Hiệp sĩ khẽ liếc nhìn về phía Mio rồi thì thầm với bạn mình: ''Họ lại bắt đầu diễn kịch nữa rồi kìa\dots''

Rushia không buồn ngẩng đầu lên, chỉ khẽ nhếch môi cười: ''Để họ cãi nhau đi-nanodesu. Cờ lật thú vị hơn nhiều.''

\ct

Hóa ra, dù miệng thì hô hào làm phim, nhưng sự thật là chẳng ai trong số họ có nổi một ý tưởng kịch bản nào ra hồn. Một cuộc họp khẩn cấp được triệu tập ngay giữa phòng khách văn phòng. Fubuki, Matsuri và Mio ngồi quây lại thành một vòng tròn, bầu không khí bỗng chốc trở nên nghiêm trọng đến mức kỳ lạ.

Nàng Lễ hội, với bản tính thích làm trung tâm, là người nổ phát súng đầu tiên: ''Tôi biết rồi! Một bộ phim điện ảnh thành công thì phải có những yếu tố mà khán giả luôn thèm khát!''

Fubuki chống cằm, ánh mắt nheo lại đầy vẻ suy tư như đang nghiền ngẫm một kế hoạch vĩ đại. Cô gật đầu tỏ vẻ đồng tình nhưng vẫn không quên thắc mắc: ''Ví dụ như?''

Mio đề xuất một cách đầy hy vọng: ''Hay là chúng ta làm về một cuộc hành trình khám phá những vùng đất mới? Kiểu như\dots\ những thứ mà chúng ta thường không làm ấy?''

Fubuki lập tức đập tay xuống sàn, reo lên: ''Đúng rồi! Vậy thì chúng ta làm một tập về\dots\ ĐỒ BƠI nhé!''

Mio lập tức dội ngay một gáo nước lạnh: ''Không được! Chúng ta đã làm cái đó quá nhiều lần rồi! Từ cái hồi dựng bể bơi trong nhà\footnote{Xem lại {\flaregothic ÉPISODE 14}, quyển số Một} cho đến mấy cái MV mùa hè\footnote{Xem lại {\flaregothic ÉPISODE 18}, quyển số Hai}, khán giả xem đến phát ngấy rồi bà ơi! Ai mà thèm xem thêm nữa chứ?''

Matsuri không chịu thua kém, chen ngang với vẻ mặt đầy mơ mộng và có phần\dots\ hơi thiên vị bản thân: ''Hay là làm về việc tất cả mọi người đều phải phục tùng tôi như một vị công chúa tối cao? Tưởng tượng cảnh Matsuri-sama ngồi trong lâu đài lộng lẫy, còn mấy người phải quỳ dưới chân tôi\dots''

Mio và Fubuki hoàn toàn phớt lờ nàng Cổ động, tiếp tục quay sang thảo luận với nhau như thể cô không hề tồn tại.

\img{3-9c}

''Vậy thì\dots'' Fubuki tiếp tục, nụ cười trên môi trở nên nghịch ngợm. ''Làm một cuộc thi sức bền xuyên biên giới thì sao?''

Mio hoang mang tột độ: ''Ngược lại chỗ nào vậy hả bà nội!? Sức bền thì liên quan gì đến phim điện ảnh?''

Bị ngó lơ, Matsuri giậm chân bành bạch như trẻ con: ''Này! Đừng có làm ngơ tôi như thế chứ! Ý tưởng công chúa của tôi tuyệt vời thế mà!''

Sự hỗn loạn càng tăng thêm khi Ayame và Choco bước vào phòng. Nàng Quỷ sừng nhìn hai người bạn đang đỏ mặt tía tai đấu khẩu, liền hỏi với giọng ngây thơ: ''Có chuyện gì mà mọi người cãi nhau như mổ bò thế? Cho tôi góp vui với nè!''

Nàng Y tá mỉm cười đầy quyến rũ, ghé sát vào vòng tròn: ''Cô cũng muốn tham gia đóng góp cho nghệ thuật nữa. Nghe có vẻ thú vị đấy chứ?''

Ayame, với vẻ mặt đầy tự tin và luôn là người nghĩ ra những ý tưởng không ai ngờ tới, giơ tay phát biểu: ''Tôi có kế hoạch này đỉnh lắm. Chúng ta sẽ làm một pha\dots\ chiếm đóng đài truyền hình quốc gia! Sau đó chúng ta sẽ bắt cả thế giới phải xem livestream của chúng ta 24/7!''

\img{3-9d}

Cả căn phòng im bặt. Mio gần như muốn ngất xỉu, cô hoảng hốt: ``Quỷ ơi, bình tĩnh lại đi! Ta không thể làm thế được! Ta là thần tượng mà, có phải khủng bố đâu!?''

Nàng Oni nhún vai, như thể đó chỉ là chuyện nhỏ: ``Sao không? Nếu muốn được cả thế giới biết đến, chúng ta phải tiếp cận tới ti-vi đầu tiên chứ, nè?''

``Ý tưởng của bà thì tôi hiểu,'' nàng Sói đen xoa thái dương, cố gắng giữ giọng bình tĩnh nhất có thể, ``nhưng như thế thì quá quắt lắm rồi!''

Nàng Y tá cười nhẹ, tiến lên phía trước, tự tin phát biểu: ``Cô hiểu mọi chuyện rồi. Để cô thử đề xuất một cái gì đó xem sao.'' Mio liếc nhìn nàng Y tá với ánh mắt nghi ngờ: ``Có được không vậy?''

Cô không đáp, chỉ mỉm cười đầy bí ẩn. Đằng sau cô, Pekora đang bận rộn gặm cà rốt, trong khi A-chan mệt mỏi gục xuống bàn vì phải tăng ca, còn Aqua vẫn chăm chỉ lau sàn như thường lệ.

Fubuki, vẫn chưa từ bỏ ý tưởng đồ bơi, quả quyết thêm lần nữa: ``Tôi vẫn muốn có một tập về đồ bơi! Đó là cách tuyệt vời để thu hút khán giả mà!''

Choco-sensei đẩy gọng kính, đưa ra một nhận định đầy chuyên môn: ''Cô nghĩ chúng ta nên giữ bộ phim ở mức nhãn \(\mathbb{PG}\)\footnote{Các hạng phim được xếp loại của các tổ chức duyệt phim (thường là MPAA - hiệp hội điện ảnh Hoa Kỳ) nhằm phần loại phim theo mức độ phổ rộng theo đối tượng khán giả đại chúng. \(\mathbb{PG}\) là dành cho mọi đối tượng nhưng có một số phân đoạn nên cân nhắc.} thôi. Đừng có dọa khán giả chạy mất dép ngay từ trailer chứ.''

Mio nghe lời giải thích của Choco-sensei liền gật gù, ánh mắt ánh lên sự ngưỡng mộ: ``Nghe cũng hợp lý đấy\dots''

Cô tiếp tục, lần này giọng đầy kiên quyết: ``Nếu chúng ta muốn thay đổi, trước hết hãy bắt đầu với tiêu đề đã.'' Matsuri nghiêng đầu hỏi: ``Như là?''

Bỗng, ánh mắt Choco tối sầm lại, giọng nói đầy khí thế: ''Chúng ta cần một cái gì đó mang tính biểu tượng hơn. Ví dụ như tiêu đề: \textit{hololive Darkness}!''

Mio hét lên: ''Cái đó chắc chắn sẽ bị xếp hạng \(\mathbb{R}\)\footnote{Tương tự như chú thích trước, với \(\mathbb{R}\) là dành cho người 17 tuổi (hay 18 tuổi ở châu Á) trở lên.} cho xem! Mà này, chẳng phải cái logo kia lại nhái từ T*-Love R* 2 sao!?''

\ct

Cuộc tranh luận càng lúc càng đi vào ngõ cụt. Giữa lúc căng thẳng nhất, cặp đôi Okayu và Korone bất ngờ xuất hiện từ phía sau ghế sofa, làm Mio giật mình suýt rơi khỏi chỗ ngồi.

Nàng Mèo vẫn giữ vẻ mặt điềm nhiên như không: ``Bọn tôi nghe hết rồi\dots'' Nàng Khuyển tiếp lời, giọng tràn đầy năng lượng: ``\dots cứ để tụi này giúp cho!''

Dù nàng Sói đen đã quen với những bất ngờ nhưng vẫn không giấu nổi sự ngạc nhiên. Cô hỏi: ``Hai em này tự dưng từ đâu ra vậy?''

Cả hai người cùng thế hệ chỉ nhếch mép cười, khoanh tay đứng đó như những nhân vật phản diện vừa lộ diện. Đằng sau họ, bộ ba Mel, Subaru, và Shion đang mải mê trò chuyện với cái loa Lolexa của nàng Phù thủy\footnote{Xem lại {\flaregothic ÉPISODE 29}, quyển số Hai.}, chẳng thèm để ý đến sự hỗn loạn phía trước.

\img{3-9f}

Okayu hắng giọng, giương lên khuôn mặt nghiêm nghị một cách bất thường: ``Mọi người không hiểu rồi. Ý tưởng của mọi người đơn giản quá.''

Korone gật gù, đồng tình một cách đầy tâm huyết: ``Đúng đó! Mọi người phải nghĩ thật xa vào! Kiểu như là\dots'' Nhưng ngay khi cô vừa mới mở miệng: ``\textit{hololive T-}'' thì Mio lập tức cắt giọng: ``Đừng nói với chị là đề xuất\dots\ tiệc trà nữa nhé?''

Bị từ chối ngay từ trong trứng nước một cách phũ phàng, cặp Chó-Mèo lập tức xụ mặt xuống, thút thít dắt tay nhau đi ra phía cửa. Okayu thở dài: ''Về thôi Koro-san'' Korone u sầu: ``Thế giới này không ai hiểu được tầm nhìn nghệ thuật của chúng mình cả\dots''

\img{3-9g}

Mio vội vàng gọi với theo trong sự hối lỗi: ''T-Tôi xin lỗi! Quay lại đi mà! Chỉ là tôi thấy nó không hợp làm phim hành động thôi!''

Nhưng hai người họ đã hoàn toàn đắm chìm vào thế giới riêng, bắt đầu bàn bạc xem nên chọn trà xanh hay trà lúa mạch cho buổi tiệc sắp tới. 

``Okayu, bà thích trà lúa mạch hay trà sấy hơn?''
  
``Đương nhiên là trà lúa mạch rồi.''

Họ dường như quá mải mê với ý tưởng, bỏ ngoài tai luôn mọi thứ xung quanh. Mio bó tay, hét lên trong sự bất lực: ``Hai người làm ơn góp ý giùm đi!''

Trong lúc đó, Matsuri vẫn cố gắng vớt vát: ''Vậy là ý tưởng công chúa của tôi được chọn rồi đúng không? Kiểu, mọi người sẽ coi tôi như là công--''

Fubuki thản nhiên quay sang hỏi, cắt ngang: ''Ủa, Matsuri-chan, bà vẫn còn ở đây à? Tôi cứ tưởng bà về từ lúc nào rồi chứ.''

Nàng Lễ hội đứng hình, hai mắt trợn ngược vì bị tổn thương sâu sắc: ``Nãy giờ tôi ở đây mà! Từ nãy đến giờ ta nói chuyện mà không để ý à!?''

Bên cạnh đó, nàng Cáo trắng vẫn không ngừng bám lấy ý tưởng cũ, quay sang người bạn của cô với ánh mắt đầy quyết tâm: ``Tôi nói rồi, chúng ta cứ làm một tập về đồ tắm đi! Khán giả sẽ thích mà!''

Nàng Sói đen, rõ ràng đã cạn kiệt sức lực, lắc đầu và than thở: ``Và tôi đã bảo chúng ta đã làm nhiều về nó rồi! Với cả, ai lại muốn mặc đồ bơi vào mùa đông cơ chứ!?''

Ayame lại chen vào: ``Vậy thì chiếm đài truyền hình cũng được ha!'' Choco vẫn giữ vẻ mặt nghiêm túc, đưa ra một ý kiến khác: ``Cô nghĩ là chúng ta nên theo đuổi ý tưởng về \textit{hololive 100\%}.''

Mio gần như hét lên: ``Ta không thể làm về chúng được! Với lại, từ từ đã, cô lại đổi tiêu đề phim nữa sao?''

Suốt một phút sau đó, cả thảy bảy người đều đưa ra ý kiến của mình, chỉ để rồi hoặc là bị nàng Sói đen phủ quyết, hoặc bị một ai khác chen ngang, đôi lúc còn vòng vo lại về đúng chủ đề mà họ định nói. 

Khi cuộc tranh cãi giữa bảy người bắt đầu lan sang cả những chuyện lông gà vỏ tỏi, cửa phòng bỗng mở ra. Sora lặng lẽ bước vào. Đôi mắt nàng Thần tượng ánh lên vẻ mệt mỏi, như thể cô vừa phải gánh vác cả bầu trời trên vai.

Thấy tiền bối xuất hiện, cả căn phòng lập tức im phăng phắc. Mio ngập ngừng lên tiếng: ''Sora-senpai\dots\ chị đến rồi sao?''

\begin{figure}[H]
  \centering
  \begin{subfigure}[t]{0.45\textwidth}
    \centering
    \includegraphics[width=8cm]{../../../../Image/Book3/3-9h1.png}
  \end{subfigure}
  \quad
  \begin{subfigure}[t]{0.45\textwidth}
    \centering
    \includegraphics[width=8cm]{../../../../Image/Book3/3-9h2.png}
  \end{subfigure}
\end{figure}

Sora khẽ gật đầu, giọng nói nhỏ nhẹ nhưng phảng phất vẻ bất lực: ``Ừ-Ừ. Chị vừa mới tới.'' Nhìn dáng vẻ thẫn thờ của Sora, Matsuri bỗng nảy ra một ý kiến: ``Sao chúng ta không để chị ấy quyết định nhỉ? Chị ấy là tiền bối mà!''

Nàng Thần tượng khẽ nhíu mày, có vẻ hơi bất ngờ trước đề xuất của Matsuri. Tuy nhiên, cô vẫn giữ giọng nhẹ nhàng, dù có vẻ đã không còn tươi tỉnh: ``Cho bộ phim à? Chị đã nghe hết ý tưởng của các em rồi\dots\ Mọi người thực sự rất tâm huyết nhỉ.''

Fubuki hồ hởi: ''Vâng ạ! Nhưng mà\dots'' rồi cô ngừng lại. Cô chợt nhớ ra điều gì đó, quay sang Miko - người đang mải mê nghịch tóc của Aki\footnote{Xem lại {\flaregothic ÉPISODE 20}, quyển số Hai.}: ``Này, Miko-chi, nãy giờ chúng ta bàn nhiều thế, nhưng bộ phim này sẽ dài trong bao lâu nhỉ?''

Miko lúc này đang mải mê nghịch tóc của Aki, cũng ngước lên tò mò: ''Chắc là phải dài hơn mấy cái clip trên kênh của mình nhiều chứ ne?''

Mio cũng tán thành với Fubuki, đồng thời quay sang Sora: ``Đúng đó. Em cũng không biết bộ phim này sẽ dài trong bao lâu nữa. Chị có biết không, Sora-senpai?'' Từ phía sau, Haato đang chật vật khiêng con cá ngừ khổng lồ\footnote{Xem lại {\flaregothic ÉPISODE 32} trong quyển này.}, giương đôi mắt cầu cứu tới mọi người.

Nàng Thần tượng lặng lẽ gật đầu. Cô nhìn mọi người bằng ánh mắt vừa mệt mỏi đến mức không còn sức sống vừa bất lực, rồi nhẹ nhàng đáp: ``Có thể là như vậy đấy.''

Fubuki háo hức đoán: ``Hai tiếng ạ? Hay ba tiếng!?''

Sora thở dài, đôi vai cô trĩu xuống. Cô nhìn mọi người bằng ánh mắt đầy thương cảm rồi thào thào: ''Hai phút rưỡi.''

Cả căn phòng, bao gồm cả những người đang lau sàn hay bưng cá ở phía sau, đồng loạt hét lên: ''\textbf{CÁI GÌ CƠ!?}''

Nhưng đó chưa phải là tất cả. Sora nhìn vào chiếc đồng hồ trên tay rồi nói tiếp: ''Và nó\dots\ cũng sắp hết rồi đấy.''

\img{3-9i}

Matsuri há hốc mồm: ``Đùa à!? Hai phút rưỡi thì chỉ đủ chiếu cái trailer thôi chứ phim ảnh gì! Nãy giờ chúng ta cãi nhau công cốc sao!?''

Mio chưa kịp phản ứng thì màn hình lớn trong phòng bỗng vụt sáng, dòng chữ credits bắt đầu chạy ngược lên phía trên cùng với những hình ảnh kỷ niệm của các thành viên.

\img{3-9j}

\begin{center}
  ''Chúng mình hy vọng các bạn sẽ tiếp tục ủng hộ \hll!\\
  Hẹn gặp lại các bạn trong những dự án tiếp theo nhé!''
\end{center}

Mọi người nhìn nhau ngơ ngác, định dọn đồ ra về thì đột nhiên, không gian chìm vào bóng tối hoàn toàn. ``Ủa? thế là hết rồi sao?'' Matsuri ngạc nhiên thốt lên, khuôn mặt vẫn còn ngơ ngác. Nhưng rồi, trong căn phòng tối om, một giọng nói vang lên. Không phải giọng của bất kỳ ai trong số họ, mà là một âm thanh lạ lùng, ấm áp và đầy uy quyền từ hệ thống loa:

\begin{center}
  ''\textbf{Chưa đâu. Đây mới chỉ là khởi đầu cho một chương hoàn toàn mới thôi!}''
\end{center}

Ánh sáng mờ ảo dần hiện lên, chiếu rọi những hình ảnh chưa từng thấy: một chiếc bàn gỗ lớn phủ khăn đỏ, trên bàn là những đĩa bánh chưng, bánh tét được sắp xếp cẩn thận. Những phong bao lì xì đỏ rực được xếp ngay ngắn bên cạnh chén trà, cùng hình ảnh những cành đào, cành mai bung nở.

Có những thứ trong số này họ đã từng thấy qua, thậm chí còn được trải nghiệm trực tiếp với chúng, như bao lì xì đỏ rực và chén trà, nhưng có những thứ mà họ chưa thực sự được nhìn thấy bao giờ, kể từ khi họ đến đây làm việc tại \hll, hay thậm chí được chiêm ngưỡng trong đời.

``Đây là\dots?'' Fubuki ngơ ngác nhìn màn hình, giọng nói xen lẫn sự ngạc nhiên và tò mò. Đến bản thân Sora cũng không giấu nổi sự bối rối: ``Ai đã quay lại đoạn này?''

Không một ai trong phòng có câu trả lời. Thậm chí cả Sora hay A-chan - những người luôn là người nắm rõ mọi thứ - cũng chỉ biết đứng đó, lặng nhìn hình ảnh trên màn hình chuyển dần sang một khung cảnh khác: những ngọn đèn lồng lung linh treo cao trong đêm tối, từng tiếng pháo hoa vang dội làm rực sáng cả bầu trời.

``Đừng nói với tôi là chúng ta lại bị quay lén nữa nha!'' Mio la lớn, ánh mắt toát vẻ lên một sự hoang mang nhẹ, nhưng cùng lúc đó, cô cảm thấy có gì đó\dots\ phấn khích. ``Thật sự thì chuyện gì đang xảy ra vậy?''

Nàng Cáo trắng khẽ lên tiếng, phá tan bầu không khí tĩnh lặng: ``Vậy\dots\ bộ phim này chưa kết thúc, đúng không?''

Sora khẽ mỉm cười, đôi mắt cô đã lấy lại được sự rạng rỡ: ''Có vẻ như\dots\ kịch bản thực sự của chúng ta giờ mới bắt đầu lộ diện đấy.''

Mio bước lên trước, ánh mắt quyết đoán hơn bao giờ hết: ``Nếu đã vậy, chúng ta phải làm cho cái kết thật xứng đáng! Mọi người, sẵn sàng chưa?''

Cả nhóm đồng thanh đáp lại, không còn chút do dự nào. Bầu không khí lúc này không còn là hỗn loạn như ban đầu, mà tràn đầy sự hứng khởi, đoàn kết và chờ đợi một điều lớn lao sắp tới.

Ở góc phòng, chiếc đồng hồ nhỏ treo trên tường bất chợt đổ chuông, báo hiệu những ngày cuối cùng của năm Kỷ Hợi Âm lịch đang trôi đi, từng giây, từng giây một\dots

%==========================================TRUYỆN PHỤ=========================================
\omk

\begin{center}
  \pov{Hikawa Kuon}
  Quận \ruby{Kyuukami}{Câu Giấy}, \ruby{Kawachi}{Hà Nội}.
\end{center}

Với nhà Hikawa chúng tôi, tiền đúng là rất quan trọng - càng nhiều thì càng vui, nhưng không có gì quan trọng hơn gia đình cả. Dẫu cho mỗi người đều có cuộc sống và công việc riêng, cứ đến dịp Tết, tất cả chúng tôi luôn trở về nhà, bất kể khoảng cách hay trở ngại.

Ở Etsunan, Tết không chỉ là một ngày lễ, mà còn là biểu tượng của sự đoàn tụ và khởi đầu mới. Bất cứ gia đình nào trong vùng cũng đều giữ thói quen này, dù nghèo khó hay dư dả. Như bài hát nổi tiếng của Bờ-lắc\footnote{Ở đây đang nói đến bài \textit{Đi về nhà} của Đen Vâu và JustaTee.} từng viết:

\hyperlink{https://www.youtube.com/watch?v=vTJdVE_gjI0}{
  \tabs[6cm] ``Đường về nhà là vào tim ta \\
  \tabs[6cm] Dẫu nắng mưa gần xa \\
  \tabs[6cm] Thất bát, vang danh, \\
  \tabs[6cm] Nhà vẫn luôn chờ ta.''}

Tôi về đến nhà vào đúng ngày 23 tháng Chạp. Nhờ cái 'cánh cửa thần kỳ' mà tôi đã dày công nghiên cứu (kết hợp một chút phép thuật mượn từ Shion), tôi có thể di chuyển giữa hai quốc gia trong nháy mắt. Ryuuhou, cô em gái nhỏ của tôi, dù vẫn phải học nhưng vì có 'ông anh làm to' ở \hll\ nên được trường đặc cách cho học trực tuyến để chuẩn bị đón Tết sớm.

Bố mẹ tôi đã bắt đầu dọn dẹp nhà cửa từ mấy hôm trước, chuẩn bị lễ cúng Ông Công ông Táo và dọn dẹp nhà cửa để đón năm mới. Đặc biệt là ông giám đốc Tanigo-san (YAGOO) của tôi, không biết ông ấy tìm hiểu văn hóa Etsunan từ bao giờ mà hào phóng cho tôi nghỉ tận một tháng.

''Cậu cứ về đón Tết cho đàng hoàng đi, Hikawa-san. Mấy cô gái ở đây cứ để tôi và A-chan lo,'' ông ấy đã nói vậy với nụ cười hiền từ quen thuộc. ``Lương của cậu vẫn được trả đủ.'' Tôi còn cảm động đến mức suýt rơi nước mắt, và cứ như vậy an tâm trở về nhà đón Tết.

Hôm nay đã là ngày 29 Tết, tính theo Âm lịch. Trong không khí tất bật của những ngày cuối năm, cả nhà tôi vẫn đang hoàn thiện những công việc cuối cùng: dọn dẹp nhà cửa, sắp xếp mâm ngũ quả và chuẩn bị bánh chưng, bánh tét. Dù mệt nhưng ai cũng cảm thấy vui vẻ, bởi đây là khoảnh khắc mà gia đình được gần gũi bên nhau nhất trong năm. Nhà cửa vẫn còn một mớ việc chưa xong, nhưng công việc thì không chờ đợi ai bao giờ.

Tôi đang hì hục lau nốt cái sàn nhà tầng hai thì điện thoại trong túi quần rung lên bần bật. Là A-san. Tôi vội đeo tai nghe vào để vừa làm vừa nói chuyện.

''Nhà Hikawa xin nghe ạ. Có chuyện gì mà chị gọi em vào giờ này thế, A-san?'' tôi mở lời bằng giọng điệu hơi mệt mỏi.

''Chào Kuon-san. Ở bên đó không khí Tết thế nào rồi?'' A-san hỏi thăm xã giao một chút, nhưng tôi thừa biết chị ấy chẳng bao giờ gọi điện chỉ để hỏi thăm sức khỏe.

''Cũng ổn chị ạ. Có chuyện gì hệ trọng sao?''

''Chà, nói sao nhỉ\dots\ Các thành viên vừa hoàn thành bộ phim điện ảnh đầu tiên của họ,'' A-san ngập ngừng.

Tôi suýt nữa thì trượt chân ngã: ''Cái gì cơ!? Làm phim á? Sao em chẳng nghe ai nói gì cả?''

''Thì lúc cậu vừa đi là họ bắt tay vào làm luôn mà. Nhưng\dots\ kết quả có vẻ không được như mong đợi. Mio-san nói rằng nó quá nhạt nhẽo, chẳng khác gì cuộc sống thường ngày.''

Tôi bật cười lớn: ''Thì đúng rồi còn gì nữa! Văn phòng mình gần nhưngày nào chẳng xảy ra mấy chuyện điên rồ, giờ có đưa lên phim thì khán giả cũng thấy bình thường thôi. Thế rồi sao nữa ạ?''

Giọng A-san bỗng trở nên nghiêm túc hơn: ''Chúng tôi cần làm một bộ phim mới. Một bộ phim thực sự có chiều sâu. Và mọi người đều nhất trí rằng\dots\ cậu là nhân vật không thể thiếu.''

Tôi sững người, chiếc chổi lau nhà rơi phịch xuống sàn: ''Hả!? Đừng có bảo là chị muốn em bay về Nhật ngay bây giờ nhé? Hôm nay mới có 29 thôi đấy!''

''Không, không cần phải về đây đâu,'' chị trấn an, nhưng câu tiếp theo của chị ấy mới thực sự là một cú nổ: ''Bởi vì tất cả chúng tôi\dots\ sẽ sang Etsunan đón Tết cùng cậu!''

''\textbf{Cái gì!?}'' tôi hét lớn đến mức Ryuuhou từ phòng bên cạnh chạy sang, nhìn tôi bằng ánh mắt lo lắng.

''Onii-san? Anh bị sao thế? Có con gián nào à?'' Ryuuhou hỏi, tay vẫn còn cầm chiếc máy tính bảng.

Tôi chưa kịp trả lời em gái thì giọng trầm ấm của YAGOO bất ngờ vang lên trong điện thoại: ''Chào Hikawa-san. Tôi đã nghe về phong tục Tết của nước cậu, và tôi nghĩ đó là một đề tài tuyệt vời cho bộ phim mới của \hll. Chúng tôi sẽ đóng quân tại nhà cậu để quay phim nhé!''

Tôi cảm thấy đầu óc mình quay cuồng: ''Nhưng mà Tanigo-san\dots\ nhà cháu nhỏ lắm, làm sao chứa hết cả mấy chục con người được!?''

''Cậu cứ khéo lo,'' YAGOO cười hiền: ''Tôi đã thuê luôn cái khu nhà bên cạnh quán cà phê sát vách nhà cậu rồi. Còn việc di chuyển thì Rushia-chan đã chuẩn bị sẵn một vòng tròn ma thuật khổng lồ. Chốc nữa chúng tôi sẽ có mặt!''

Tôi đứng chết lặng. A-san bồi thêm một câu trước khi cúp máy: ''Chuẩn bị tiếp khách đi nhé, Kuon-san.''

*Tút\dots\ tút\dots\ tút\dots*

Tôi nhìn chiếc điện thoại với ánh mắt vô hồn. Cô em gái bước tới, tò mò nhìn màn hình: ''Sao thế anh? Có chuyện gì vui à?''

Tôi thở dài, quay sang nhìn cô em gái: ''Ryuuhou này\dots\ Em chuẩn bị tinh thần đi. Lúc nữa nhà mình sẽ trở thành phim trường\dots\ và chúng ta sẽ phải tiếp đón một đám người 'không bình thường' nhất thế giới đấy.''

Mắt Ryuuhou bỗng sáng rực lên: ''Hả? Sora-chan và mọi người sắp đến thật sao!? Tuyệt quá! Để em đi bảo mẹ chuẩn bị thêm bánh chưng!''

Nhìn cô em gái hớn hở chạy xuống lầu, tôi chỉ biết tựa lưng vào tường mà than thở: ''Trời ạ\dots\ Đến cả nghỉ Tết ở quê nhà mà cũng không thoát khỏi họ nữa!''

Nhưng trong lòng tôi, một cảm giác háo hức lạ kỳ bỗng trỗi dậy. Tết năm nay, chắc chắn sẽ là cái Tết náo nhiệt và điên rồ nhất lịch sử nhà Hikawa cho mà xem.

Muốn biết chuyện gì đã xảy ra khi đám người \hll\ đổ bộ xuống ở Kawachi quê tôi? Hẹn gặp lại các bạn trong quyển T - ''Ngày Tết quê Kuon'' nhé!

% \il

%===========================================TRAILER===========================================
% \begin{center}
%     \flaregothic\fontsize{20pt}{20pt}\selectfont \color{eptype}PROCHAINEMENT DANS...\\
%     \gangofthree\fontsize{60pt}{60pt}\selectfont \color{epnum}QUYỂN T\color{black}
% \end{center}

% \vspace{10pt}

% \txtbx[3]{
%     \centering
%     \textit{''Khi những bước chân đầu tiên chạm vào mảnh đất thủ đô...''}
% }

% \vspace{15pt}

% \begin{itemize}
%     \item \textbf{[Cảnh quay nhanh: Những con phố Hà Nội tấp nập xe máy, tiếng còi xe inh ỏi xen lẫn không khí se lạnh của những ngày cuối năm.]}
%     \item \textbf{[Cánh cửa thần kỳ bật mở. Hơn hai mươi cô gái ngã nhào lên nhau giữa phòng khách nhà Hikawa.]}
%     \item \textbf{Noel \& Roboco:} ''Ááá! Cái chậu nước này ở đâu ra vậy nè!?''
%     \item \textbf{Kuon:} ''Trời đất ơi! Cái sàn nhà tôi vừa mới lau xong!!!''
% \end{itemize}

% \vspace{10pt}

% \txtbx[2]{
%     \centering \vnmsans \fontsize{16pt}{16pt}\selectfont
%     \textbf{MỘT TRẢI NGHIỆM TẾT ĐỘC NHẤT VÔ NHỊ}
% }

% \vspace{10pt}

% \begin{itemize}
%     \item \textbf{[Matsuri, Pekora và Shion đứng hình trước căn phòng 'siêu nhỏ' của cậu Tổng.]}
%     \item \textbf{Shion:} ''Nhà anh... trông có vẻ 'tiết kiệm' diện tích quá nhỉ?''
%     \item \textbf{Kuon:} ''Này, đừng có dùng từ 'nghèo' một cách tế nhị như thế chứ!''
% \end{itemize}

% \vspace{10pt}

% \begin{itemize}
%     \item \textbf{[Cảnh quay các nhóm idol tất bật dọn dẹp: Okayu \& Korone sơn tường, Mio \& Choco lau bếp, Miko \& Mel lau bàn thờ.]}
%     \item \textbf{[Cả gia đình quây quần bên nồi bánh chưng nghi ngút khói.]}
%     \item \textbf{Họ hàng nhà Kuon:} ''Thế... các em đã có gấu chưa?''
%     \item \textbf{Sora, Mel, Mio:} ''[Mặt đỏ bừng, lắp bắp] D-Dạ... chưa ạ!''
% \end{itemize}

% \vspace{15pt}

% \txtbx[3]{
%     \centering
%     \textit{''Nhưng giữa những tiếng cười rộn rã...\\...luôn có những nốt trầm lặng lẽ.''}
% }

% \vspace{10pt}

% \begin{itemize}
%     \item \textbf{[Âm thanh pháo hoa nổ tung xen lẫn tiếng hét của Pekora.]}
%     \item \textbf{[Màn hình mờ dần đi, chỉ còn lại âm thanh tít tút của máy đo nhịp tim.]}
%     \item \textbf{Kuon (giọng nghẹn ngào):} ''Có những cái Tết... sẽ không bao giờ giống như cũ nữa.''
% \end{itemize}

% \vspace{20pt}

% \begin{center}
%     \gangofthree\fontsize{48pt}{48pt}\selectfont \color{epnum}MỪNG TẾT CANH TÝ\color{black}\\
%     \vnmsans\fontsize{18pt}{18pt}\selectfont (Extra Volume - Quyển T)\\
%     \continuum\fontsize{14pt}{14pt}\selectfont Đón đọc ngay sau Book 3!
% \end{center}

% \vspace{20pt}

\end{document}
